% ============================================================
% APPLICATION PACKAGE ENGINE
% Read-only for job-specific generators. Reusable rendering commands only.
% ============================================================

% Shared bilingual command \MyL is defined in ../helper/common_commands.tex.

% ------------------------------------------------------------
% WRAPPER PAGE NUMBERS
% The page counter is still used internally for TOC destinations, but no page
% number is printed when OnePDF.GeneratedPageNumbers.Visible = 0.
% This does not alter page numbers already embedded in source PDFs.
% ------------------------------------------------------------
\ifnum\Flag{OnePDF.GeneratedPageNumbers.Visible}=1\relax
  \pagestyle{plain}
\else
  \pagestyle{empty}
\fi

% ------------------------------------------------------------
% COMBINED COVER PAGE + TABLE OF CONTENTS
% The separate TOC switch in ../src/content_onePDF.tex stays false because the
% TOC is rendered directly on this first page. Compile twice after structural
% or document-selection changes.
% ------------------------------------------------------------
\newcommand{\ApplicationCoverPage}{%
  \begin{titlepage}
    \thispagestyle{empty}
    \centering
    \vspace*{1.0cm}

    {\Huge\textbf{\ApplicantFullName}\par}
    \vspace{0.75cm}

    {\Large\textbf{\MyL{Application}{Bewerbung}}\par}
    \vspace{0.25cm}

    {\Large
      \ApplicationPosition\par
    }

    \vspace{0.55cm}

    {\large
      \ApplicationOrganisation\par
      \ApplicationDepartment\par
    }

    \IfStrEq{\ApplicationReference}{}{}{%
      \vspace{0.25cm}
      {\normalsize
        \MyL{Reference}{Referenz}: \ApplicationReference\par
      }
    }

    \vspace{0.85cm}

    \begin{minipage}{0.90\textwidth}
      \raggedright
      \tableofcontents
    \end{minipage}

    % \tableofcontents may switch the current page to the plain page style.
    % Force the combined cover/TOC page back to empty so no wrapper page
    % number is printed.
    \thispagestyle{empty}

    \vfill

    {\normalsize
      \ApplicationLocation\,\textperiodcentered\,\today\par
    }
  \end{titlepage}%
}

% ------------------------------------------------------------
% GENERIC PDF INSERTION
% A document is included only if the file exists.
% ------------------------------------------------------------
\newcommand{\IncludeApplicationDocument}[2]{%
  \IfFileExists{#2}{%
    \includepdf[
      pages=-,
      addtotoc={1,section,1,{#1},appdoc:#2}
    ]{#2}%
  }{%
    \PackageWarning{application}{Document '#2' not found -- skipped}%
  }%
}

\newcommand{\IncludeSubDocument}[3]{%
  \IfFileExists{#3}{%
    \includepdf[
      pages=-,
      addtotoc={1,#1,2,{#2},subdoc:#3}
    ]{#3}%
  }{%
    \PackageWarning{application}{Document '#3' not found -- skipped}%
  }%
}

% ------------------------------------------------------------
% ACADEMIC QUALIFICATIONS
%
% Document hierarchy and selection are declared in application_documents.tex (project registry; selection comes from ../src/content_onePDF.tex).
%
% The first existing academic PDF creates the section entry.
% No extra separator pages are inserted.
%
% LLM EDITING RULE:
% - Complete-section visibility is controlled in ../src/content_onePDF.tex.
% - Individual academic document visibility is controlled in ../src/content_onePDF.tex.
% - Do not delete standard file definitions or reusable macros.
% ------------------------------------------------------------
\newif\ifAcademicSectionStarted
\AcademicSectionStartedfalse

\NewDocumentCommand{\AcademicDocument}{m m O{-}}{%
  % #1 = TOC title
  % #2 = PDF file
  % #3 = OPTIONAL pdfpages page selection, e.g. [1-3,8-11]
  %      Omit it to include the complete PDF.
  %
  % Selected page ranges should start with source page 1 so the TOC anchor remains stable.
  \IfFileExists{#2}{%
    \ifAcademicSectionStarted
      \includepdf[
        pages={#3},
        addtotoc={1,subsection,2,{#1},academicitem:\thepage}
      ]{#2}%
    \else
      \global\AcademicSectionStartedtrue
      \includepdf[
        pages={#3},
        addtotoc={
          1,section,1,{\MyL{\Value{OnePDF.Section.AcademicQualifications.TitleEN}}{\Value{OnePDF.Section.AcademicQualifications.TitleDE}}},academic,
          1,subsection,2,{#1},academicitem:\thepage
        }
      ]{#2}%
    \fi
  }{%
    \PackageWarning{application}{Academic document '#2' not found -- skipped}%
  }%
}

\newcommand{\AcademicQualifications}{%
  \AcademicQualificationDocuments
}

% ------------------------------------------------------------
% PROFESSIONAL EXPERIENCE EVIDENCE
%
% Data model:
%   Employer Certificate
%     -> certificate issued by the Contractual employer
%     -> may cover several project assignments
%
%   Project Certificate
%     -> optional additional certificate issued for one
%        Operational company / project assignment
%
% Examples:
%   one contractual-employer certificate may support several role assignments
%   assignment-specific project certificates remain optional additions
% ------------------------------------------------------------
\newif\ifProfessionalSectionStarted
\ProfessionalSectionStartedfalse

% Internal helper. The first existing professional PDF creates the
% Professional Experience section entry; later PDFs become subsections.
\newcommand{\ProfessionalEvidenceDocument}[2]{%
  \IfFileExists{#2}{%
    \ifProfessionalSectionStarted
      \includepdf[
        pages=-,
        addtotoc={1,subsection,2,{#1},professional:#2}
      ]{#2}%
    \else
      \global\ProfessionalSectionStartedtrue
      \includepdf[
        pages=-,
        addtotoc={
          1,section,1,{\MyL{\Value{OnePDF.Section.ProfessionalExperience.TitleEN}}{\Value{OnePDF.Section.ProfessionalExperience.TitleDE}}},professional,
          1,subsection,2,{#1},professional:#2
        }
      ]{#2}%
    \fi
  }{%
    \PackageWarning{application}{Professional evidence '#2' not found -- skipped}%
  }%
}

% Employer Certificate
% #1 Contractual employer
% #2 PDF file
\newcommand{\EmployerCertificate}[2]{%
  \ProfessionalEvidenceDocument
    {\MyL{Employer Certificate}{Arbeitszeugnis} -- #1}%
    {#2}%
}

% Optional Project Certificate
% #1 Contractual employer
% #2 Operational company / project assignment
% #3 PDF file
\newcommand{\ProjectCertificate}[3]{%
  \ProfessionalEvidenceDocument
    {\MyL{Project Certificate}{Projektzeugnis} -- #1 -- \MyL{assignment at}{Einsatz bei} #2}%
    {#3}%
}

% Backward-compatible generic command.
% Prefer \EmployerCertificate and \ProjectCertificate in new selections.
\newcommand{\ProfessionalCertificate}[3]{%
  \IfStrEq{#1}{#2}{%
    \EmployerCertificate{#1}{#3}%
  }{%
    \ProjectCertificate{#1}{#2}{#3}%
  }%
}

% ------------------------------------------------------------
% REFEREE CONTACT INFORMATION
% Referees are added in application_documents.tex (project registry; selection comes from ../src/content_onePDF.tex) using \Referee.
%
% Layout goal:
% - shared typography and section-title style from ../helper/style.tex
% - page title at the top
% - referee blocks placed around the vertical middle of the page
% - moderate spacing between referees for easy scanning
%
% Syntax:
% \Referee
%   {Name}
%   {Organisation EN}{Organisation DE}
%   {Department EN}{Department DE}
%   {E-mail}
%   [optional status/note, already bilingual if required]
%
% LLM EDITING RULE:
% - Keep contact details exactly as provided; do not invent titles, positions,
%   departments, e-mail addresses, or current employment status.
% - Individual referee visibility is controlled in ../src/content_onePDF.tex.
% - Contact-page visibility is controlled in ../src/content_onePDF.tex.
% ------------------------------------------------------------
\newcommand{\RefereeList}{}

\NewDocumentCommand{\Referee}{m m m m m m O{}}{%
  \gappto\RefereeList{%
    \noindent
    \begin{minipage}{\linewidth}
      {\large\textbf{-- #1}}\par
      \vspace{0.22cm}
      \MyL{#2}{#3}\par
      \MyL{#4}{#5}\par
      \ifstrempty{#7}{}{#7\par}%
      \texttt{#6}\par
    \end{minipage}
    \par\vspace{1.05cm}%
  }%
}

\newcommand{\RefereeContactPage}{%
  \clearpage
  \phantomsection
  \ifAcademicReferenceSectionStarted
    \addcontentsline{toc}{subsection}{\MyL{Referee Contact Information}{Kontaktinformationen der Referenzpersonen}}%
  \else
    \global\AcademicReferenceSectionStartedtrue
    \addcontentsline{toc}{section}{\MyL{\Value{OnePDF.Section.AcademicReferences.TitleEN}}{\Value{OnePDF.Section.AcademicReferences.TitleDE}}}%
    \addcontentsline{toc}{subsection}{\MyL{Referee Contact Information}{Kontaktinformationen der Referenzpersonen}}%
  \fi

  % Page title follows the general application/CV typography.
  \section*{\MyL{\Value{OnePDF.Section.AcademicReferences.TitleEN}}{\Value{OnePDF.Section.AcademicReferences.TitleDE}}}
  % \vspace{0.15cm}
  % {\normalsize\MyL{Referee Contact Information}{Kontaktinformationen der Referenzpersonen}\par}

  % Keep the contact blocks visually centered on the sheet instead of
  % accumulating directly below the title.
  \vfill
  \begin{center}
    \begin{minipage}{0.78\textwidth}
      \RefereeList
    \end{minipage}
  \end{center}
  \vfill

  \clearpage
}

% ------------------------------------------------------------
% ACADEMIC REFERENCE LETTERS
% #1 Referee name
% #2 PDF file
% ------------------------------------------------------------
\newif\ifAcademicReferenceSectionStarted
\AcademicReferenceSectionStartedfalse

\newcommand{\AcademicReferenceLetter}[2]{%
  \IfFileExists{#2}{%
    \ifAcademicReferenceSectionStarted
      \includepdf[
        pages=-,
        addtotoc={1,subsection,2,{\MyL{Academic Reference Letter}{Akademisches Empfehlungsschreiben} -- #1},reference:#2}
      ]{#2}%
    \else
      \global\AcademicReferenceSectionStartedtrue
      \includepdf[
        pages=-,
        addtotoc={
          1,section,1,{\MyL{\Value{OnePDF.Section.AcademicReferences.TitleEN}}{\Value{OnePDF.Section.AcademicReferences.TitleDE}}},references,
          1,subsection,2,{\MyL{Academic Reference Letter}{Akademisches Empfehlungsschreiben} -- #1},reference:#2
        }
      ]{#2}%
    \fi
  }{%
    \PackageWarning{application}{Academic reference letter '#2' not found -- skipped}%
  }%
}

% ------------------------------------------------------------
% ADDITIONAL CERTIFICATES
%
% These are supplementary qualifications/awards which are not part of the
% core academic graduation documents and are not employment certificates.
% Each physical certificate is uploaded as a separate PDF and registered in
% application_documents.tex (project registry; selection comes from ../src/content_onePDF.tex).
%
% Prepared files are registered with generic numeric IDs in application_documents.tex (project registry; selection comes from ../src/content_onePDF.tex).
%
% LLM EDITING RULE:
% - Section and document visibility are controlled in ../src/content_onePDF.tex.
% - application_documents.tex contains only project-local registration/rendering logic.
% - Do not delete source PDFs simply because they are not used in one job.
% ------------------------------------------------------------
\newif\ifAdditionalCertificateSectionStarted
\AdditionalCertificateSectionStartedfalse

% #1 English title
% #2 German title
% #3 PDF file
\NewDocumentCommand{\AdditionalCertificate}{m m m o}{%
  % #1 = English title
  % #2 = German title
  % #3 = PDF file path
  % #4 = OPTIONAL: single PDF page number
  %
  % If #4 is omitted, all pages of the PDF are included.
  % Examples:
  %   \AdditionalCertificate{English}{Deutsch}{file.pdf}
  %       -> all pages
  %
  %   \AdditionalCertificate{English}{Deutsch}{file.pdf}[2]
  %       -> only page 2
  %
  % LLM NOTE:
  % Leave the optional page argument [<page>] away when the complete
  % certificate should be included. Add [<page>] only if a single
  % relevant page should be extracted from a multi-page source PDF.

  \IfNoValueTF{#4}{%
    \def\AdditionalCertificatePages{-}%
    \def\AdditionalCertificateTOCPage{1}%
  }{%
    \def\AdditionalCertificatePages{#4}%
    \def\AdditionalCertificateTOCPage{#4}%
  }%

  \IfFileExists{#3}{%
    \ifAdditionalCertificateSectionStarted
      \includepdf[
        pages={\AdditionalCertificatePages},
        addtotoc={
          \AdditionalCertificateTOCPage,
          subsection,
          2,
          {\MyL{#1}{#2}},
          additionalcertificate:#3
        }
      ]{#3}%
    \else
      \global\AdditionalCertificateSectionStartedtrue
      \includepdf[
        pages={\AdditionalCertificatePages},
        addtotoc={
          \AdditionalCertificateTOCPage,
          section,
          1,
          {\MyL{\Value{OnePDF.Section.AdditionalCertificates.TitleEN}}{\Value{OnePDF.Section.AdditionalCertificates.TitleDE}}},
          additionalcertificates,
          \AdditionalCertificateTOCPage,
          subsection,
          2,
          {\MyL{#1}{#2}},
          additionalcertificate:#3
        }
      ]{#3}%
    \fi
  }{%
    \PackageWarning{application}
      {Additional certificate '#3' not found -- skipped}%
  }%
}

% ============================================================
% SECTION RENDERING / ORDER CONTROL
% ============================================================

\newcommand{\RenderCoverLetter}{%
  \IfFlag{OnePDF.Document.CoverLetter.Visible}{%
    \IncludeApplicationDocument
      {\MyL{Cover Letter}{Anschreiben}}
      {\Value{OnePDF.File.CoverLetter}}%
  }{}%
}

\newcommand{\RenderCV}{%
  \IfFlag{OnePDF.Document.CV.Visible}{%
    \IncludeApplicationDocument
      {\MyL{Curriculum Vitae}{Lebenslauf}}
      {\Value{OnePDF.File.CV}}%
  }{}%
}

\newcommand{\RenderAcademicQualifications}{%
  \IfFlag{OnePDF.Section.AcademicQualifications.Visible}{\AcademicQualifications}{}%
}

\newcommand{\RenderProfessionalExperience}{%
  \IfFlag{OnePDF.Section.ProfessionalExperience.Visible}{\ProfessionalDocuments}{}%
}

\newcommand{\RenderAcademicReferences}{%
  \IfFlag{OnePDF.Section.AcademicReferences.Visible}{%
    \IfFlag{OnePDF.RefereeContactInformation.Visible}{\RefereeContactPage}{}%
    \AcademicReferenceDocuments
  }{}%
}

\newcommand{\RenderAdditionalCertificates}{%
  \IfFlag{OnePDF.Section.AdditionalCertificates.Visible}{\AdditionalCertificateDocuments}{}%
}

% Dispatch one ordered item ID from ../src/content_onePDF.tex.
\newcommand{\RenderOnePDFItem}[1]{%
  \ifcsname RenderOnePDFItem@#1\endcsname
    \csname RenderOnePDFItem@#1\endcsname
  \else
    \PackageWarning{application}{Unknown onePDF order item '#1' -- skipped}%
  \fi
}

\expandafter\def\csname RenderOnePDFItem@CoverLetter\endcsname{\RenderCoverLetter}
\expandafter\def\csname RenderOnePDFItem@CV\endcsname{\RenderCV}
\expandafter\def\csname RenderOnePDFItem@AcademicQualifications\endcsname{\RenderAcademicQualifications}
\expandafter\def\csname RenderOnePDFItem@ProfessionalExperience\endcsname{\RenderProfessionalExperience}
\expandafter\def\csname RenderOnePDFItem@AcademicReferences\endcsname{\RenderAcademicReferences}
\expandafter\def\csname RenderOnePDFItem@AdditionalCertificates\endcsname{\RenderAdditionalCertificates}

\newcommand{\RenderOnePDFPackage}{%
  \RenderOnePDFItem{\Value{OnePDF.Order01}}%
  \RenderOnePDFItem{\Value{OnePDF.Order02}}%
  \RenderOnePDFItem{\Value{OnePDF.Order03}}%
  \RenderOnePDFItem{\Value{OnePDF.Order04}}%
  \RenderOnePDFItem{\Value{OnePDF.Order05}}%
  \RenderOnePDFItem{\Value{OnePDF.Order06}}%
}
